% TeX root=../main.tex

O sistema do eslavo antigo evoluiu ao curso dos três séculos de sua história e
não é monolítico em razão das divergências dialetais.
No entanto não há necessidade de separar radicalmente um sistema original,
aquele de Cirilo e Metódio, de um sistema mais tardio que será aquele que
atestam os textos conservados.

\section{Consonantismo}

\subsection{Inventário de fonemas}

\begin{table}[!h]
  \caption{Sistema de consoantes}\label{tab:consoantes}
  \begin{center}
    \begin{tabular}[c]{cccccc}
      \toprule
       & Labiais & \multicolumn{2}{c}{Dentais}  & Palatais & Velares \\
       &  & não palatalizadas  & palatalizadas & \multicolumn{1}{c}{} &  \\
       \midrule
       oclusivas & p & t & & & k\\
        & b & d & & & g\\
        \midrule
       fricativas & f (?) & s & (s') & š & x\\
        & v & z & & ž & \\
        \midrule
       africadas &  & c &  & č & \\
        &  & dz & &  & \\
        \midrule
       nasais & m & n & n' &  & \\
        \midrule
       líquidas &  & r & r' &  & \\
        &  & l & l' &  & \\
        \midrule
       semivogais &  &  &  & j & \\
       \bottomrule
    \end{tabular}
  \end{center}
\end{table}

O sistema de consoantes do eslavo antigo compo rta 22 ou 24 fonemas segundo os
autores, que os dividem em quatro pontos de articulação (labial, dental,
palatal e velar) e dois modos de articulação (surdo e sonoro).\footnote{%
  A apresentação do consonantismo que figura no \emph{Manuel} de
  \textcite[57]{Vaillant1964} é falsa do ponto de vista fonológico. Não há uma
  oposição sistemática entre consoantes duras, moles e palatalizadas (p.~60).
}
A oposição entre surda e sonora vale por todas as consoantes salvo às
ressoantes e as consoantes africada /č/ e a fricativa /x/, que não possuem
correspondente sonoro.
Para as consoantes da série dental há uma oposição entre fonemas palatalizados
e não palatalizados (/r/ \emph{vs} /r'/, /l/ \emph{vs} /l'/, /n/ \emph{vs}
/n'/: \textsc{gen.sg.masc.} \cirilex{polu}\slash{}\ocsex{polu} `da metade'
\emph{vs} \textsc{dat.sg.neu.} \cirilex{pol'ju}\slash{}\ocsex{pol'u} `no
campo', \textsc{acu.sg.fem.}
\cirilex{materI}\slash{}\ocsex{materI} `mãe' \emph{vs}
adj.~\textsc{nom.sg.masc.}
\cirilex{mater'I}\slash{}\ocsex{mater'I} `materno',
\textsc{gen.sg.masc.} \cirilex{gospodina}\slash{}\ocsex{gospodina} `do
mestre' \emph{vs} adj.~\textsc{gen.sg.masc}
\cirilex{gospodin'ja}\slash{}\ocsex{gospodin'a} `de mestre').

O estatuto de /j/ ainda é objeto de discussões: em eslavo comum ele não é um
fonema mas um alofone de /ĭ/, que continua a alofonia herdade do indo-europeu
entre *\emph{i} e *\pietrans{y}, mas em eslavo antigo o fonema /ĭ/ não existe
mais, ele se tornou \slash{}\ocstrans{I}/ e a questão da autonomia de /j/ se impõe.
É de fato impossível opor /\ocstrans{I}/ e /j/ em um mesmo contexto, isto é,
formando um par mínimo, o que argumenta contra o estatuto de fonema.
Mas para o eslavo antigo a maioria dos autores admitem que /j/ é um fonema pleno.
O estatuto de fonema de /j/ é assegurado pela queda dos \emph{jers} que leva
ao enfraquecimento do /\ocstrans{I}/ fraco e a confusão do /\ocstrans{I}/
forte com /e/: à partir do momento em que não há mais o /\ocstrans{I}/, não há
mais a alofonia, e /j/, mesmo admitindo que ele não teria sido até então nada
se não um alofone de /\ocstrans{I}/, adquire desse fato o estatuto de fonema.
Antes da queda dos \emph{jers}, é mais difícil de provar o estatuto fonemático
de /j/.
Em todo caso, no século XI, /j/ é um fonema pleno.

O estatuto de /f/ é discutível.
Ele não aparece se não em empréstimos do grego e não é conhecido na língua
vernácula.
A transcrição do /f/ grego é também variável e, se na língua culta, ele é
representado por /f/, grafado \grafema{\cirilex{f}}, incluindo no alfabeto
glagolítico que comporta um grafema específico, na língua popular ele é
representado pela oclusiva labial correspondente /p/ (\emph{cf.} em russo antigo
o antropônimo \ocsex{*stepanU} de \emph{Stephanos}).
No entanto, os nomes próprios como \emph{Sofia} (Santa Sofia de Constantinopla)
ou os empréstimos como \cirilex{fevrar'I}\slash{}\ocsex{fevrar'I} `fevereiro'
aclimataram rapidamente o /f/ e pode-se considerar que desde antes do período
eslavo antigo o /f/ fazia sim parte dos fonemas da língua, ao menos no socioleto
culto.

O eslavo comum tardio tinha um fonema /s'/, produzido a partir de *\emph{x} pela
terceira palatalização, que encontramos notavelmente no pronome
\cirilex{vIsI}\slash{}\ocsex{vIsI} `todo'.
Esse fonema /s'/ não existia em eslavão morávio, pois nele o resultado de *\emph{x}
após a terceira palatalização é /š/.
Em eslavo meridional, em contraste, esse fonema existia, mas foi eliminado.
Para o \textsc{nom.sg.fem.} /v\ocstrans{I}s'a/, o alfabeto glagolítico, que não
diferencia entre /ě/ e /ja/, grafa \ocsex{vIsE}; o alfabeto cirílico, que faz
diferença entre /ě/ e /ja/ e possui dois grafemas separados, possui duas
grafias, \cirilex{vIsE}\slash{}\ocsex{vIsE}, que segue a grafia glagolítica, mas que é
morfologicamente irregular (não existem \textsc{nom.sg.fem.} em -\ocsex{E}), ou
\cirilex{vIsa}\slash{}\ocsex{vIsa}, morfologicamente regular, mas que não representa a
palatalização de /s'/; não encontramos em eslavo antigo a grafia
*\cirilex{vIsja}\slash{}\ocsex{vIs'a}, que só aparece em eslavão.
Da mesma maneira, para o \textsc{acu.sg.fem.}, encontramos as duas grafias
\cirilex{vIsjõ}\slash{}\ocsex{vIs'õ} e \cirilex{vIsõ}\slash{}\ocsex{vIsõ}, a última não
representando a palatalização de /s'/.
Encontramos igualmente as grafias \ocsex{vIsa} e \ocsex{vIsõ} nos manuscritos
glagolíticos e o mesmo para o adjetivo \ocsex{vIs'akU} `cada um, todo
(indefinido)' grafado também \ocsex{vIsEkU} ou \ocsex{vIsakU}.
As grafias \ocsex{vIsa} e \ocsex{vIsõ} indicam que o /s'/ foi despalatalizado e
confundido com /s/: o sistema fonológico comportaria um /s'/ no primeiro período
do eslavo antigo, mas não mais ao final do período eslavo antigo (é por isso que
ele aparece em parêntese na~\autoref{tab:consoantes}).

O eslavo antigo de Cirilo e Metódio possuiria provavelmente também (no século
IX) um /t'/ e um /d'/ palatalizados.
O tratamento de \ocsex{St}, \ocsex{Zd}, com um grupo bifonemático,
característico do búlgaro e não do macedônio, é mais tardio: ele é sistemático
em todos os textos eslavos antigos (dos quais nenhum é anterior ao fim do século
X).

As velares /k/ e /g/ possuem os alofones palatalizados [k'] e [g'] antes de
vogais anteriores, a palatalização sendo por vezes explicitamente representada na
grafia com o diacrítico \grafema{{\rus{}  ҄ }}.
As velares não se encontram nesses contextos se não em palavras emprestadas e
[k'] e [g'] não são se não variantes de /k/ e /g/ e não fonemas da língua.

As líquidas /r/ e /l/ podem funcionar como centro de sílaba e tem os alofones
vocálicos \pietrans{R} e \pietrans{L} e, para as líquidas palatalizadas,
\pietrans{R'} e \pietrans{L'}.
A questão das líquidas vocálicas é há muito disputada, mas a maioria dos
pesquisadores hoje aceitam que o eslavo antigo possuíam em verdade líquidas
vocálicas, como o checo ou o servo-croata modernos.
É necessário no entanto ressaltar que se tratam de alofones e não fonemas plenos:
não existe uma oposição entre /r/ e */\emph{\pietrans{R}}/ e entre /l/ e
*/\emph{\pietrans{L}}/, que estão em distribuição complementar: encontra-se
\pietrans{R} e \pietrans{L} apenas entre duas consoantes, nos outros contextos
encontra-se \emph{r} e \emph{l}.
Esses alofones vocálicos eram realizados na articulação da líquida seguida de
uma vogal reduzida, daí as grafias \grafema{\ocsex{rU}}, \grafema{\ocsex{rI}},
\grafema{\ocsex{lU}} e \grafema{\ocsex{lI}} (\cirilex{srIdIce}\slash{}\ocsex{srIdIce}
`coração' < \glsxtrshort{eslavocomum} *\emph{sĭrdĭce} < *\pietrans{cRdi-ko-},
\cirilex{vlIkU}\slash{}\ocsex{vlIkU} `lobo' < \glsxtrshort{eslavocomum} *\emph{vĭlko-}
< *\pietrans{wLko-}).
Nada distingue na grafia as líquidas de sequência etimológica /r/ +
/\ocstrans{U}/ ou /\ocstrans{I}/ e /l/ + /\ocstrans{U}/ ou /\ocstrans{I}/,
mas sabe-se pela evolução posterior que a confusão era apenas gráfica e não
fonética, dado que o ``\emph{jer}'' das antigas líquidas vocálicas não entra na
alternância normal dos \emph{jers} fortes e fracos (cf.~\ref{sec:jers}), ao
contrário de um \emph{jer} de uma sequência etimológica /r/ + /\ocstrans{U}/,
/r/ + /\ocstrans{I}/:  o \emph{jer} de \cirilex{krUvI}\slash{}\ocsex{krUvI} `sangue', que possui um
*\emph{rŭ} etimológico, é vocalizado em \cirilex{krovI}\slash{}\ocsex{krovI} desde o
eslavo antigo, enquanto o ``\emph{jer}'' de \cirilex{mrItvU}, \cirilex{mrUtvU}\slash{}\ocsex{mrItvU}, \ocsex{mrUtvU} `morto', que possui um \pietrans{R} etimológico,
não o é jamais (sobre a variação de
\grafema{\ocstrans{rI}}/\grafema{\ocstrans{rU}}, \emph{cf.} acima).

O inventário não é exatamente o mesmo nos dialetos.
Assim, /dz/ é simplificado a /z/ em búlgaro antigo (os grafemas
\grafema{\cirilex{ꙁ}}, \grafema{\cirilex{dz}} são desconhecidos de \glsxtrshort{Sav}
e \glsxtrshort{Supr}).
Ele está mais bem conservado em macedônio antigo, onde há uma flutuação entre a
africada /dz/ e a contínua /z/ e alguns casos de hipercorreção (/dz/ no lugar de
/z/ etimológico), mas essa flutuação em si sugere que a distinção entre os dois
fonemas estava em vias de desaparecer, com prováveis divergências regionais.
Em eslavo antigo morávio, o \glsxtrlong{Kiev} não possui /z/.

Uma outra divergência diz respeito à oposição entre /r/ e /r'/, que tende a
desaparecer em búlgaro antigo em favor de /r/ (grafia
\cirilex{cEsara}\slash{}\ocsex{cEsara} no lugar de \cirilex{cEsar'ja}\slash{}\ocsex{cEsar'a}
`César', \glsxtrshort{Supr}) ao passo que ela é conservada nos textos macedônios antigos.
Essa neutralização vale também para o alofone vocálico:
a grafia \grafema{\ocstrans{rU}} é bem mais frequente que a grafia
\grafema{\ocstrans{rI}}, e nesse aspecto o macedônio antigo está de acordo com o
búlgaro antigo; assim encontra-se mais frequentemente
\cirilex{srUdIce}\slash{}\ocsex{srUdIce} `coração' do que
\cirilex{crIdIce}\slash{}\ocsex{srIdIce}, \cirilex{mrUtvU}\slash{}\ocsex{mrUtvU} `morto do que
\cirilex{mrItvU}\slash{}\ocsex{mrItvU} (manteremos aqui a grafia
\grafema{\ocstrans{rI}} mesmo que ela seja minoritária, seguindo a prática da
maioria dos autores).
A tendência à neutralização existe também, de uma maneira menos afirmada, para
*\pietrans{L} e *\pietrans{L'}, \grafema{\ocstrans{lU}} é notavelmente mais
frequente que \grafema{\ocstrans{lI}}, assim \cirilex{vlUkU}\slash{}\ocsex{vlUkU}
`lobo' no lugar de \cirilex{vlIkU}\slash{}\ocsex{vlIkU}.

O eslavo antigo, como outras línguas eslavas, é caracterizado por um número
importante de fonemas palatais e palatalizados.
O eslavo pertence ao grupo \emph{satəm} do indo-europeu e as palatais antigas
*\pietrans{c}, *\pietrans{j} e *\pietrans{jh} aparecem na forma de sibilantes
/s/ e /z/, como em iraniano.
Mas o eslavo já tinha desde muito a tendência à palatalização e palatalizou, em
certos contextos, como o indo-iraniano, as velares resultantes de labiovelares
indo-europeias, tal como todos os grupos de C + /j/
(\emph{cf.}~\ref{sec:consoantes}).
A consequência das palatalizações sucessivas é uma série de alternâncias
sistemáticas nas flexões nominal e verbal e na derivação.
O resultado da palatalização é, em geral, um fonema único, mas em cestos casos
resultou em um grupo bifonemático: assim a palatalização das dentais produziu em
búlgaro \ocsex{St} e \ocsex{Zt} (que é um desenvolvimento mais tardio,
\emph{cf.} acima), em contrates às africadas \emph{c} e \emph{dz} do dialeto
morávio; a palatalização das labiais resulto em um grupo bifonemático
\emph{pl'}, \emph{bl'}, \emph{vl'} e \emph{ml'} com um [l] epentético.
Notar-se-á que em certos manuscritos o [l] epentético está por vezes ausente na
fronteira de morfemas (\textsc{loc.sg} \ocsex{na zemi} `na terra', Mt 6, 10
\glsxtrshort{Marianus}, em contraste com \ocsex{na zeml'i}, \glsxtrshort{Zogr}), em particular em
\glsxtrshort{Supr}, onde um \grafema{\ocstrans{U}} marca em sequência a lenição da
labial (\textsc{1sg} \cirilex{postavIjõ}\slash{}\ocsex{postavIjõ} `instituirei'
\passagem{Supr124}, no lugar de \cirilex{postavl'jõ}\slash{}\ocsex{postavl'õ}).

\subsection{Alternâncias morfofonológica}

\subsubsection{Alternâncias devido palatalizações}
\lipsum[1]
\subsubsection{Alternâncias posicionais}
\lipsum[2]

\subsection{Distribuição}
\lipsum[4]
\subsection{Grafia}
\lipsum[1]

\section{Vocalismo}

\subsection{Inventário de fonemas}
\lipsum[3]
\subsection{Alternâncias vocálicas}
\lipsum[3]
\subsection{Distribuição}
\lipsum[5]
\subsection{Grafia}
\lipsum[1]
\subsection{Os \emph{jers}}\label{sec:jers}
\lipsum[3]

\section{Traços supra-segmentais}
\lipsum[1]

